\chapter{Экспериментальный выбор архитектуры}\label{ch:experiments}

\section{Методика}

Задание допускает произвольную конфигурацию сети при условии наличия не менее двух
скрытых слоёв, поэтому архитектура выбиралась экспериментально: сравнивались шесть
конфигураций, различающихся глубиной, шириной и функцией активации скрытых слоёв.

Условия сравнения были уравнены: разбиение выборки фиксировано --- одно и то же
стратифицированное разделение 456 к 113 при сиде 42, параметры стандартизации
вычислены по обучающей части. Все конфигурации обучались 84 эпохи порциями по 8 примеров
при $\eta = 0{,}0314$, бинарной кросс-энтропии, инициализации He и нулевых смещениях.

Результат одного запуска зависит от сида, определяющего начальные веса и порядок
перемешивания, причём разброс по сиду сопоставим с различиями между конфигурациями.
Поэтому каждая конфигурация обучалась пять раз с зёрнами 1, 7, 42, 123 и 2026, в
сравнении участвуют средние значения, а дополнительно вычислялось среднее квадратичное
отклонение функции потерь на проверке --- оно показывает, какие различия в таблице
значимы. Сравнение оформлено как автоматический тест.

\section{Результаты сравнения}

Результаты приведены в таблице~\ref{tab:architectures}. Величины $L$ и $L_{\text{val}}$
обозначают функцию потерь на обучающей и проверочной выборках соответственно, $s$~---
среднее квадратичное отклонение $L_{\text{val}}$ по пяти запускам, $A_{\text{val}}$~---
доля правильных ответов на проверочной выборке.

\begin{table}[htbp]
\caption{Сравнение архитектур, средние значения по пяти запускам}
\label{tab:architectures}
\begin{center}
\begin{tabular}{|l|r|r|r|r|r|}
\hline
Архитектура & Параметров & $L$ & $L_{\text{val}}$ & $s$ & $A_{\text{val}}$ \\
\hline
24--24, сигмоида      & 1394 & 0,0519 & 0,0792 & 0,0018 & 0,9664 \\
24--24--24, сигмоида  & 1994 & 0,0515 & 0,0821 & 0,0046 & 0,9646 \\
48--48, сигмоида      & 3938 & 0,0530 & 0,0768 & 0,0037 & 0,9646 \\
8--8, сигмоида        &  338 & 0,0522 & 0,0788 & 0,0075 & 0,9664 \\
24--24, ReLU          & 1394 & 0,0034 & 0,0674 & 0,0293 & 0,9752 \\
24--24, тангенс       & 1394 & 0,0075 & \textbf{0,0265} & 0,0097 & \textbf{0,9858} \\
\hline
\end{tabular}
\end{center}
\end{table}

\section{Выводы}

Ёмкость сети не является ограничивающим фактором. Четыре конфигурации с
сигмоидой различаются по числу параметров в двенадцать раз --- от 338 до 3938, --- однако
функция потерь на проверке меняется лишь в пределах 0,0768--0,0821, что сопоставимо со
средним квадратичным отклонением по сиду (до 0,0075) и потому незначимо. Решающей
является функция потерь на \emph{обучающей} выборке: она составляет около 0,052 у всех
четырёх конфигураций независимо от размера. Будь сеть ограничена ёмкостью,
двенадцатикратный рост числа параметров позволил бы лучше приблизить обучающую выборку.
Этого не происходит, значит ограничение лежит не в выразительной способности модели, а в
сходимости: при заданных скорости обучения и числе эпох обучение не успевает завершиться.

Добавление слоя ухудшает результат: переход от двух скрытых слоёв к трём при
той же ширине увеличивает число параметров с 1394 до 1994, но функция потерь на проверке
растёт с 0,0792 до 0,0821, а разброс по сиду --- с 0,0018 до 0,0046. Причина та же:
производная сигмоиды не превосходит $0{,}25$, поэтому каждый дополнительный слой
уменьшает проходящий градиент не менее чем вчетверо. Третий слой замедляет сходимость, не
давая взамен выигрыша в качестве.

Определяющим параметром оказалась функция активации: замена сигмоиды на
гиперболический тангенс при неизменной архитектуре снижает функцию потерь на проверке
втрое --- с 0,0792 до 0,0265, --- повышает долю правильных ответов с 0,9664 до 0,9858, а
на обучающей выборке даёт падение с 0,0519 до 0,0075, то есть сеть действительно начинает
сходиться. Объяснение следует из~(\ref{eq:activation-grad}): максимум производной
тангенса равен единице против $0{,}25$ у сигмоиды, что даёт четырёхкратный выигрыш в
величине градиента на каждом слое. Кроме того, тангенс симметричен относительно нуля,
тогда как сигмоида выдаёт только положительные значения, из-за чего компоненты градиента
по весам одного нейрона сонаправлены и спуск идёт зигзагом.

ReLU сходится быстрее, но обобщает хуже: она даёт наименьшую из всех функцию
потерь на обучающей выборке (0,0034), но на проверочной --- 0,0674: разрыв между
выборками 0,064 против 0,019 у тангенса, то есть переобучение, а разброс по сиду
($s = 0{,}0293$) втрое больше. Причина --- обнуление производной на отрицательной
полуоси: часть нейронов при неудачной инициализации перестаёт получать градиент, и итог
существенно зависит от начальных весов.

Принята конфигурация из двух скрытых слоёв по 24 нейрона с
гиперболическим тангенсом: лучшие значения обеих метрик, наименьший разрыв между
выборками среди сошедшихся конфигураций, не менее двух скрытых слоёв. Она установлена в
программе как используемая по умолчанию.
